% =============================================================================
%  s02_desarrollo.tex — Sección 2: Desarrollo / Argumentación central
%  Agrega tantos frames como necesites dentro de cada subsección.
% =============================================================================

\section{Plantilla: Marco Conceptual}


% ── Frame: Contexto ───────────────────────────────────────────────────────────
% NOTA \pause: revela bullets uno a uno durante la presentación en vivo,
%   pero genera páginas PDF duplicadas. Descomenta solo si presentas en proyector.
\begin{frame}{Contexto y Problemática}
  \begin{columns}[T]
    \begin{column}{0.58\textwidth}
      \begin{itemize}
        \item La \zmvm\ crece a una tasa anual de \textbf{1.4\%}, con el 80\% del suelo
              nuevo en zonas periféricas sin servicio de transporte masivo.
        % \pause
        \item La expansión horizontal eleva los \alert{costos de infraestructura}
              y reduce la accesibilidad para los estratos más vulnerables.
        % \pause
        \item El \termino{modelo de densificación compacta} emerge como alternativa
              con respaldo técnico y normativo creciente \citep{lore2024densidad}.
      \end{itemize}
    \end{column}
    \begin{column}{0.38\textwidth}
      \begin{block}{Dato clave}
        El 60\% de los viajes en la \zmvm\ se realizan en transporte informal,
        señal de un sistema de movilidad estructuralmente insuficiente.
      \end{block}
    \end{column}
  \end{columns}
\end{frame}


% ── Frame: Tres modelos ───────────────────────────────────────────────────────
\begin{frame}{Modelos de Densificación Evaluados}
  \begin{columns}[T]
    \begin{column}{0.32\textwidth}
      \begin{block}{Modelo 1: TOD}
        \textit{Transit-Oriented\\Development}\\[0.2cm]
        Copenhague, Estocolmo\\[0.1cm]
        {\small Principio: compacidad alrededor de nodos de transporte masivo.}
      \end{block}
    \end{column}
    \begin{column}{0.32\textwidth}
      \begin{block}{Modelo 2: Barrio}
        \textit{Rehabilitación\\incremental}\\[0.2cm]
        Medellín, Bogotá\\[0.1cm]
        {\small Principio: participación comunitaria y urbanismo social.}
      \end{block}
    \end{column}
    \begin{column}{0.32\textwidth}
      \begin{block}{Modelo 3: Corredor}
        \textit{Corredor\\urbano mixto}\\[0.2cm]
        Ciudad de México\\[0.1cm]
        {\small Principio: mixticidad de usos sobre ejes de movilidad masiva.}
      \end{block}
    \end{column}
  \end{columns}
\end{frame}

% ── Frame: Tabla comparativa ─────────────────────────────────────────────────
\begin{frame}{Comparativa de Modelos}
  \begin{table}
    \centering
    \small
    \begin{tabularx}{\textwidth}{@{} l X X X @{}}
      \toprule
      \textbf{Criterio}   & \textbf{TOD} & \textbf{Barrio} & \textbf{Corredor} \\
      \midrule
      Equidad             & Media        & Alta            & Media--Alta \\
      Eficiencia energét. & Alta         & Media           & Alta \\
      Viabilidad inst.    & Media        & Alta            & Media \\
      Escala de impacto   & Metropolit.  & Local           & Zonal \\
      \bottomrule
    \end{tabularx}
    \fuentefigura{Elaboración propia con base en \citet{lore2024densidad}.}
  \end{table}
\end{frame}

% ── Frame: Cita destacada ────────────────────────────────────────────────────
\begin{frame}{Argumento Central}
  \vfill
  \begin{center}
    \color{ColorPrincipal}
    \Large\itshape
    ``La ciudad compacta no es una utopía formal,\\
    sino una decisión política sobre cómo se distribuyen\\
    los recursos del suelo urbano.''
  \end{center}
  \vfill
  \flushright{\small\citet{lore2024densidad}}
\end{frame}

% ── Frame con figura (plantilla) ─────────────────────────────────────────────
\begin{frame}{Resultados del Análisis Espacial}
  \begin{columns}[T]
    \begin{column}{0.55\textwidth}
      % Descomenta cuando tengas la imagen:
      % \includegraphics[width=\textwidth]{img/nombre_figura.png}
      % \fuentefigura{Elaboración propia, 2026.}

      \begin{alertblock}{Placeholder de figura}
        Inserta aquí tu imagen con:\\
        \texttt{\textbackslash includegraphics\{img/figura.png\}}
      \end{alertblock}
    \end{column}
    \begin{column}{0.42\textwidth}
      \begin{itemize}
        \item Hallazgo principal 1 \citep{lore2024zmvm}
        % \pause  ← descomenta solo para presentación en vivo (genera páginas extra en PDF)
        \item Hallazgo principal 2
        % \pause
        \item Implicación para la política pública
      \end{itemize}
    \end{column}
  \end{columns}
\end{frame}

% ── Frame: Dos imágenes con texto explicativo ─────────────────────────────────
%    USO: comparativas visuales, antes/después, dos mapas, dos planos.
%    Ajusta height para que ambas imágenes queden alineadas verticalmente.
%    INSTRUCCIÓN: cuando tengas tus imágenes, comenta el alertblock y descomenta \includegraphics.
\begin{frame}{Comparativa Visual — Dos Imágenes}
  \begin{columns}[T]
    \begin{column}{0.48\textwidth}
      \centering
      % \includegraphics[width=\textwidth, height=3.8cm, keepaspectratio]{img/figura1.png}
      \begin{alertblock}{Placeholder imagen 1}
        \texttt{img/figura1.png}
      \end{alertblock}
      {\small Texto explicativo de la imagen izquierda. Describe el contexto, hallazgo o período temporal.}\\[0.05cm]
      \fuentefigura{Elaboración propia, 2026.}
    \end{column}
    \begin{column}{0.48\textwidth}
      \centering
      % \includegraphics[width=\textwidth, height=3.8cm, keepaspectratio]{img/figura2.png}
      \begin{alertblock}{Placeholder imagen 2}
        \texttt{img/figura2.png}
      \end{alertblock}
      {\small Texto explicativo de la imagen derecha. Describe el contexto, hallazgo o período temporal.}\\[0.05cm]
      \fuentefigura{Fuente secundaria, 2026.}
    \end{column}
  \end{columns}
\end{frame}


% ── Frame: Diagrama de flujo ──────────────────────────────────────────────────
%    USO: metodología, proceso de análisis, etapas de proyecto.
%    Modifica los nodos, colores y flechas según tu contenido.
%    Requiere \usepackage{tikz} + \usetikzlibrary{positioning,shapes.geometric,arrows.meta,calc}
%    (ya incluidos en Latex/BeamerTheme.tex)
\begin{frame}{Metodología — Diagrama de Flujo}
  \begin{center}
  \begin{tikzpicture}[
    % Estilos de nodos
    caja/.style   = {rectangle, rounded corners=5pt,
                     draw=ColorPrincipal, line width=0.8pt,
                     fill=ColorFondo, text width=2.6cm,
                     align=center, minimum height=0.75cm, font=\small},
    inicio/.style = {caja, fill=ColorPrincipal, text=white, font=\small\bfseries},
    result/.style = {caja, fill=ColorAcento,    text=black,  font=\small\bfseries},
    decision/.style = {diamond, draw=ColorPrincipal, line width=0.8pt,
                       fill=ColorFondo, aspect=2.5,
                       text width=2.2cm, align=center, font=\scriptsize},
    % Estilo de flechas
    flecha/.style = {-Stealth, thick, color=ColorPrincipal},
    lateral/.style= {-Stealth, thick, color=ColorAcento, dashed},
    % Separación entre nodos
    node distance = 0.55cm and 2.2cm
  ]
    % Columna principal (flujo vertical)
    \node[inicio]   (diag)  {Diagnóstico};
    \node[caja, below=of diag]    (anal)  {Análisis};
    \node[decision, below=of anal](dec)   {¿Viable?};
    \node[caja, below=of dec]     (prop)  {Propuesta};
    \node[result, below=of prop]  (res)   {Resultado};

    % Flechas principales
    \draw[flecha] (diag) -- (anal);
    \draw[flecha] (anal) -- (dec);
    \draw[flecha] (dec)  -- node[right, font=\scriptsize]{Sí} (prop);
    \draw[flecha] (prop) -- (res);

    % Rama de ajuste (decisión → No → vuelve a diagnóstico)
    \node[caja, right=of dec] (ajuste) {Ajuste\\de enfoque};
    \draw[lateral] (dec)    -- node[above, font=\scriptsize]{No} (ajuste);
    \draw[lateral] (ajuste) |- (diag);
  \end{tikzpicture}
  \end{center}
\end{frame}
